% Complete reviewed Unit 4; editorial master: chapters/ch04/answer-key.md
% Markdown SHA-256: d32b83bfe09437697c011d6f5b29eea95e0e2d11e7ab0cc303d04487390a2da7
% One-use initial assembly followed by complete content and page review.
\newcommand{\runningtitle}{Unit 4. Answer key}
% Layout inherited from the reviewed previous course; set a running title per unit.
% This is a small hand-maintained TeX layout, not a conversion pipeline.
\RequirePackage{fix-cm}
\documentclass[11pt,twoside]{article}
\usepackage{fontspec}
\setmainfont{Cambria}[
  Renderer=HarfBuzz,
  BoldFont=cambriab.ttf,
  ItalicFont=cambriai.ttf,
  BoldItalicFont=cambriaz.ttf
]
\setmonofont{lmmono10-regular.otf}[BoldFont=lmmonolt10-bold.otf]
\usepackage[paperwidth=165mm,paperheight=241mm,inner=21mm,outer=19mm,
  top=22mm,bottom=20mm,headheight=14pt,headsep=7mm,footskip=10mm]{geometry}
\usepackage{xcolor,array,booktabs,longtable,tabularx}
\usepackage{enumitem,fancyhdr}
\usepackage[unicode,hidelinks]{hyperref}
\definecolor{ink}{HTML}{000000}
\definecolor{muted}{HTML}{4C5558}
\setlength{\parindent}{0pt}
\setlength{\parskip}{3.5pt}
\AtBeginDocument{\fontsize{11}{14.3}\selectfont}
\setlength{\tabcolsep}{4pt}
\renewcommand{\arraystretch}{1.10}
\setlist[itemize]{leftmargin=1.3em,label=\textbullet,itemsep=2pt,topsep=3pt,parsep=0pt}
\setlist[enumerate]{leftmargin=1.8em,itemsep=3pt,topsep=4pt,parsep=0pt}
\widowpenalty=10000
\clubpenalty=10000
\emergencystretch=1em
\raggedbottom
\pagestyle{fancy}
\fancyhf{}
\providecommand{\runningtitle}{Old English}
\fancyhead[LE,RO]{\small\itshape \runningtitle}
\fancyhead[LO,RE]{\small Old English}
\fancyfoot[LE,RO]{\small\thepage}
\renewcommand{\headrulewidth}{0.3pt}
\newcommand{\booktitle}[2]{
  {\fontsize{9}{11}\selectfont\color{muted}#1\par}
  \vspace{3pt}
  {\fontsize{20}{24}\selectfont\color{ink}#2\par}
  \vspace{6pt}
}
\newcommand{\source}[1]{\par{\footnotesize\color{muted}#1\par}}
\newcommand{\readingtext}[1]{
  \par\addvspace{5pt}
  \begingroup\leftskip=1.5em
    \fontsize{12}{16}\selectfont #1\par
  \endgroup\addvspace{5pt}
}
\makeatletter
\renewcommand{\section}{\@startsection{section}{1}{\z@}
  {-2.4ex \@plus -0.6ex \@minus -0.2ex}{0.8ex \@plus 0.2ex}
  {\normalfont\fontsize{12}{15}\selectfont\bfseries\color{ink}}}
\renewcommand{\subsection}{\@startsection{subsection}{2}{\z@}
  {-1.8ex \@plus -0.5ex \@minus -0.2ex}{0.5ex \@plus 0.2ex}
  {\normalfont\fontsize{11}{14}\selectfont\bfseries\color{ink}}}
\makeatother

\hyphenpenalty=10000
\exhyphenpenalty=10000
\emergencystretch=3em
\newcommand{\itempara}[2]{\par\addvspace{3pt}\noindent\hangindent=1.8em\hangafter=1\makebox[1.8em][l]{#1.}#2\par}
\newcommand{\reserve}[1]{\par\penalty-100\begingroup\dimen0=\pagegoal\advance\dimen0 by-\pagetotal\ifdim\dimen0<#1\newpage\fi\endgroup}
\begin{document}
\addfontfeatures{Ligatures=TeXOff}
\raggedright
\booktitle{UNIT 4}{Answer key}

\reserve{7\baselineskip}
% mdparagraph:P001:begin
\section*{Using the key}
% mdparagraph:P001:end

% mdparagraph:P002:begin
The models use the lesson's selected forms. A documented alternative can be valid Old English while differing from an explicitly requested selected form. In particular, \textbf{findan} has selected first-singular preterite \textbf{fand} and third-singular \textbf{funde}; first-singular \textbf{funde} and third-singular \textbf{fand} are also valid. \textbf{Helpan} allows a genitive as well as a dative, although the lesson selects the dative construction. Do not confuse a missed target with impossible language.
% mdparagraph:P002:end

% mdparagraph:P003:begin
Preserve vowel length and the distinction between \textbf{i/\allowbreak{}ī}, \textbf{ie/\allowbreak{}īe} and \textbf{ea/\allowbreak{}ēa}. The tables label every requested person and number. A column headed “1/\allowbreak{}2/\allowbreak{}3pl” includes all three plural persons; dual subjects use these plural forms. Word order may vary when cases, agreement, reference and meaning remain clear. Where a task requests a limiting genitive before its head, a grammatical postposed genitive still misses that positioning request.
% mdparagraph:P003:end

% mdparagraph:P004:begin
Infinitive and past participle columns are supplied citations. \textbf{Bregdan} and \textbf{fēolan} remain recognition references. The source reading's local help remains available; neither its harder verb forms nor its source-only constructions become independent-recall requirements. The answers and translations below are newly prepared models, not Cook's own answer key.
% mdparagraph:P004:end

\reserve{7\baselineskip}
% mdparagraph:P005:begin
\section*{Principal parts and tense}
% mdparagraph:P005:end

\reserve{5\baselineskip}
\par\begin{minipage}{\linewidth}\raggedright
% mdparagraph:P006:begin
\textbf{P1.}
% mdparagraph:P006:end

% item: NEW-U04-P01
% key: KEY-NEW-U04-P01
% mdparagraph:P007:begin
The four positions are:
% mdparagraph:P007:end

% mdtable:T01:begin
\par\begingroup
\fontsize{10.5}{13.2}\selectfont
\setlength{\LTleft}{0pt}\setlength{\LTright}{\fill}
\noindent
\begin{tabular}{@{}>{\raggedright\arraybackslash}p{22.13mm}>{\raggedright\arraybackslash}p{22.13mm}>{\raggedright\arraybackslash}p{74.12mm}@{}}
\toprule
\textbf{Position} & \textbf{Form} & \textbf{Description} \\
\midrule
1 & drīfan & Infinitive citation \\
2 & drāf & Preterite singular principal part; here first or third person \\
3 & drifon & Preterite plural principal part; any of the three persons \\
4 & drifen & Past participle citation \\
\bottomrule
\end{tabular}
\par\endgroup

% mdtable:T01:end
% mdparagraph:P008:begin
\textbf{Drāf} and \textbf{drifon} are finite indicative forms. \textbf{Drīfan} and \textbf{drifen} are nonfinite citations; neither identifies a person and number in the indicative. The table does not by itself teach an infinitive clause, perfect or passive. [1]
% mdparagraph:P008:end

\end{minipage}\par\addvspace{4pt}

\par\begin{minipage}{\linewidth}\raggedright
% mdparagraph:P009:begin
\textbf{P2.}
% mdparagraph:P009:end

% item: NEW-U04-P02
% key: KEY-NEW-U04-P02
% mdtable:T02:begin
\par\begingroup
\fontsize{10.5}{13.2}\selectfont
\setlength{\LTleft}{0pt}\setlength{\LTright}{\fill}
\noindent
\begin{tabular}{@{}>{\raggedright\arraybackslash}p{24.05mm}>{\raggedright\arraybackslash}p{24.05mm}>{\raggedright\arraybackslash}p{38.96mm}>{\raggedright\arraybackslash}p{28.5mm}@{}}
\toprule
\textbf{Verb} & \textbf{Class} & \textbf{Second singular preterite} & \textbf{Stem used} \\
\midrule
rīdan & I & ride & rid-, as in ridon \\
bēodan & II & bude & bud-, as in budon \\
bindan & III & bunde & bund-, as in bundon \\
\bottomrule
\end{tabular}
\par\endgroup

% mdtable:T02:end
% mdparagraph:P010:begin
Each takes the plural-preterite stem plus \textbf{-e}. The singular principal parts \textbf{rād, bēad, band} do not supply the requested second-person stems. [1]
% mdparagraph:P010:end

\end{minipage}\par\addvspace{4pt}

\par\begin{minipage}{\linewidth}\raggedright
% mdparagraph:P011:begin
\textbf{P3.}
% mdparagraph:P011:end

% item: NEW-U04-P03
% key: KEY-NEW-U04-P03
% mdparagraph:P012:begin
\textbf{Drāf} is first- or third-person singular preterite indicative. \textbf{Drifon} is first-, second- or third-person plural preterite indicative. A dual subject also takes the plural verb form; there is no separate dual ending.
% mdparagraph:P012:end

% mdparagraph:P013:begin
\textbf{Drifen} is the supplied past participle. Its \textbf{-en} is not another personal ending of the finite paradigm. Do not treat the resemblance to \textbf{drifon} as identity of function.
% mdparagraph:P013:end

\end{minipage}\par\addvspace{4pt}

\par\begin{minipage}{\linewidth}\raggedright
% mdparagraph:P014:begin
\textbf{P4.}
% mdparagraph:P014:end

% item: NEW-U04-P04
% key: KEY-NEW-U04-P04
% mdparagraph:P015:begin
Both verbs are second-person singular preterite indicative: “you rode” and “you heard.” \textbf{Ride} has the short \textbf{i} of \textbf{ridon}, plus \textbf{-e}. The strong verb changes its stem vowel and adds no dental preterite marker. \textbf{Hīerdest} is the weak preterite of \textbf{hīeran}, with the dental \textbf{-d-} and second-singular \textbf{-est}.
% mdparagraph:P015:end

% mdparagraph:P016:begin
Quantity matters: preterite \textbf{ride} has short \textbf{i}, unlike present \textbf{rīde}.
% mdparagraph:P016:end

\end{minipage}\par\addvspace{4pt}

\reserve{7\baselineskip}
% mdparagraph:P017:begin
\section*{Class I}
% mdparagraph:P017:end

\reserve{11\baselineskip}
\par\begin{minipage}{\linewidth}\raggedright
% mdparagraph:P018:begin
\textbf{A1.}
% mdparagraph:P018:end

% item: NEW-U04-A01
% key: KEY-NEW-U04-A01
% mdtable:T03:begin
% course-table-style: 2026-10
% course-panel:T03-P01:begin
\par\begingroup
\fontsize{10.5}{13.2}\selectfont
\setlength{\LTleft}{0pt}\setlength{\LTright}{\fill}
\noindent
\begin{tabular}{@{}>{\raggedright\arraybackslash}p{39.459mm}>{\raggedright\arraybackslash}p{39.459mm}>{\raggedright\arraybackslash}p{39.459mm}@{}}
\toprule
\textbf{Person and number} & \textbf{Present} & \textbf{Preterite} \\
\midrule
1 sg. & drīfe & drāf \\
2 sg. & drīfst & drife \\
3 sg. & drīfð & drāf \\
1–3 pl. & drīfað & drifon \\
\bottomrule
\end{tabular}
\par\endgroup
% course-panel:T03-P01:end
% mdtable:T03:end
% mdparagraph:P019:begin
Present \textbf{ī} contrasts with preterite singular \textbf{ā} and plural-based \textbf{i}. The second singular is \textbf{drife}, not a form made from \textbf{drāf}. [2]
% mdparagraph:P019:end

\end{minipage}\par\addvspace{4pt}

\par\begin{minipage}{\linewidth}\raggedright
% mdparagraph:P020:begin
\textbf{A2.}
% mdparagraph:P020:end

% item: NEW-U04-A02
% key: KEY-NEW-U04-A02
% mdtable:T04:begin
\par\begingroup
\fontsize{10.5}{13.2}\selectfont
\setlength{\LTleft}{0pt}\setlength{\LTright}{\fill}
\noindent
\begin{tabular}{@{}>{\raggedright\arraybackslash}p{14.44mm}>{\raggedright\arraybackslash}p{14.44mm}>{\raggedright\arraybackslash}p{14.44mm}>{\raggedright\arraybackslash}p{14.44mm}>{\raggedright\arraybackslash}p{15.59mm}>{\raggedright\arraybackslash}p{15.59mm}>{\raggedright\arraybackslash}p{18.19mm}@{}}
\toprule
\textbf{Verb} & \textbf{Present 1sg} & \textbf{Present 2sg} & \textbf{Present 3sg} & \textbf{Preterite 1sg} & \textbf{Preterite 2sg} & \textbf{Preterite 1/\allowbreak{}2/\allowbreak{}3pl} \\
\midrule
rīdan & rīde & rītst & rītt & rād & ride & ridon \\
wrītan & wrīte & wrītst & wrītt & wrāt & write & writon \\
\bottomrule
\end{tabular}
\par\endgroup

% mdtable:T04:end
% mdparagraph:P021:begin
Stem-final \textbf{d} with \textbf{-st} gives \textbf{-tst}; \textbf{t} with \textbf{-st} likewise gives \textbf{-tst}. Before third-singular \textbf{-ð}, the dental combination gives selected \textbf{-tt}. The documented third-singular spellings \textbf{rīt} and \textbf{wrīt} are also valid, though these models retain the selected double \textbf{t}. [2]
% mdparagraph:P021:end

\end{minipage}\par\addvspace{4pt}

\par\begin{minipage}{\linewidth}\raggedright
% mdparagraph:P022:begin
\textbf{A3.}
% mdparagraph:P022:end

% item: NEW-U04-A03
% key: KEY-NEW-U04-A03
% mdtable:T05:begin
\par\begingroup
\fontsize{10.5}{13.2}\selectfont
\setlength{\LTleft}{0pt}\setlength{\LTright}{\fill}
\noindent
\begin{tabular}{@{}>{\raggedright\arraybackslash}p{21.11mm}>{\raggedright\arraybackslash}p{21.11mm}>{\raggedright\arraybackslash}p{21.11mm}>{\raggedright\arraybackslash}p{22.8mm}>{\raggedright\arraybackslash}p{26.6mm}@{}}
\toprule
\textbf{Verb} & \textbf{Present 2sg} & \textbf{Present 3sg} & \textbf{Preterite 2sg} & \textbf{Preterite 1/\allowbreak{}2/\allowbreak{}3pl} \\
\midrule
rīsan & rīst & rīst & rise & rison \\
snīðan & snīst & snīð & snide & snidon \\
\bottomrule
\end{tabular}
\par\endgroup

% mdtable:T05:end
% mdparagraph:P023:begin
In \textbf{snīðan}, \textbf{ð} changes to \textbf{d} in the plural-preterite stem and therefore in \textbf{snide} and \textbf{snidon}. The first/\allowbreak{}third singular \textbf{snāð} retains \textbf{ð}. This lexical alternation is distinct from the present ending combinations. \textbf{Rīsan} retains \textbf{s} in the forms shown. Documented third-singular \textbf{rīsð} and \textbf{snīðð} are valid alternatives to the selected forms. [2]
% mdparagraph:P023:end

\end{minipage}\par\addvspace{4pt}

\par\begin{minipage}{\linewidth}\raggedright
% mdparagraph:P024:begin
\textbf{A4.}
% mdparagraph:P024:end

% item: NEW-U04-A04
% key: KEY-NEW-U04-A04
% mdtable:T06:begin
% course-table-style: 2026-10
% course-panel:T06-P01:begin
\par\begingroup
\fontsize{10.5}{13.2}\selectfont
\setlength{\LTleft}{0pt}\setlength{\LTright}{\fill}
\noindent
\begin{tabular}{@{}>{\raggedright\arraybackslash}p{39.459mm}>{\raggedright\arraybackslash}p{39.459mm}>{\raggedright\arraybackslash}p{39.459mm}@{}}
\toprule
\textbf{Person and number} & \textbf{Present} & \textbf{Preterite} \\
\midrule
1 sg. & wrēo & wrāh \\
2 sg. & wrīhst & wrige \\
3 sg. & wrīhð & wrāh \\
1–3 pl. & wrēoð & wrigon \\
\bottomrule
\end{tabular}
\par\endgroup
% course-panel:T06-P01:end
% mdtable:T06:end
% mdparagraph:P025:begin
The contracted first singular and plural are \textbf{wrēo} and \textbf{wrēoð}; do not append the ordinary \textbf{-e/\allowbreak{}-að} to them. Present \textbf{wrīhst/\allowbreak{}wrīhð}, preterite singular \textbf{wrāh}, and plural-based \textbf{wrige/\allowbreak{}wrigon} preserve their different vowels and \textbf{h/\allowbreak{}g} stems. [2]
% mdparagraph:P025:end

\end{minipage}\par\addvspace{4pt}

\par\begin{minipage}{\linewidth}\raggedright
% mdparagraph:P026:begin
\textbf{A5.}
% mdparagraph:P026:end

% item: NEW-U04-A05
% key: KEY-NEW-U04-A05
% mdtable:T07:begin
\par\begingroup
\fontsize{10.5}{13.2}\selectfont
\setlength{\LTleft}{0pt}\setlength{\LTright}{\fill}
\noindent
\begin{tabular}{@{}>{\raggedright\arraybackslash}p{45.23mm}>{\raggedright\arraybackslash}p{27.5mm}>{\raggedright\arraybackslash}p{45.65mm}@{}}
\toprule
\textbf{Meaning} & \textbf{Preterite} & \textbf{Corresponding present} \\
\midrule
I rode /\allowbreak{} ride & Ic rād. & Ic rīde. \\
You, one person, rode /\allowbreak{} ride & Ðū ride. & Ðū rītst. \\
He rode /\allowbreak{} rides & Hē rād. & Hē rītt. \\
We rode /\allowbreak{} ride & Wē ridon. & Wē rīdað. \\
\bottomrule
\end{tabular}
\par\endgroup

% mdtable:T07:end
% mdparagraph:P027:begin
The second singular past uses \textbf{ride}, not \textbf{rād}. Present \textbf{rīt} is a valid alternative to selected third-singular \textbf{rītt}. The pronouns distinguish first and third singular where the past form coincides.
% mdparagraph:P027:end

\end{minipage}\par\addvspace{4pt}

\par\begin{minipage}{\linewidth}\raggedright
% mdparagraph:P028:begin
\textbf{A6.}
% mdparagraph:P028:end

% item: NEW-U04-A06
% key: KEY-NEW-U04-A06
% mdtable:T08:begin
\par\begingroup
\fontsize{10.5}{13.2}\selectfont
\setlength{\LTleft}{0pt}\setlength{\LTright}{\fill}
\noindent
\begin{tabular}{@{}>{\raggedright\arraybackslash}p{39.46mm}>{\raggedright\arraybackslash}p{39.46mm}>{\raggedright\arraybackslash}p{39.46mm}@{}}
\toprule
\textbf{Subject} & \textbf{Present} & \textbf{Preterite} \\
\midrule
ic & drīfe & drāf \\
wē & drīfað & drifon \\
ðū & drīfst & drife \\
hē & drīfð & drāf \\
\bottomrule
\end{tabular}
\par\endgroup

% mdtable:T08:end
% mdparagraph:P029:begin
Keep \textbf{wē} plural in both tenses. The change to the preterite does not change the subject's person or number, and the second singular requires the plural-based stem.
% mdparagraph:P029:end

\end{minipage}\par\addvspace{4pt}

\reserve{7\baselineskip}
% mdparagraph:P030:begin
\section*{Class II}
% mdparagraph:P030:end

\reserve{11\baselineskip}
\par\begin{minipage}{\linewidth}\raggedright
% mdparagraph:P031:begin
\textbf{B1.}
% mdparagraph:P031:end

% item: NEW-U04-B01
% key: KEY-NEW-U04-B01
% mdtable:T09:begin
% course-table-style: 2026-10
% course-panel:T09-P01:begin
\par\begingroup
\fontsize{10.5}{13.2}\selectfont
\setlength{\LTleft}{0pt}\setlength{\LTright}{\fill}
\noindent
\begin{tabular}{@{}>{\raggedright\arraybackslash}p{39.459mm}>{\raggedright\arraybackslash}p{39.459mm}>{\raggedright\arraybackslash}p{39.459mm}@{}}
\toprule
\textbf{Person and number} & \textbf{Present} & \textbf{Preterite} \\
\midrule
1 sg. & bēode & bēad \\
2 sg. & bīetst & bude \\
3 sg. & bīett & bēad \\
1–3 pl. & bēodað & budon \\
\bottomrule
\end{tabular}
\par\endgroup
% course-panel:T09-P01:end
% mdtable:T09:end
% mdparagraph:P032:begin
The present second and third singular, \textbf{bīetst/\allowbreak{}bīett}, have \textbf{īe} where the infinitive citation has \textbf{ēo}. The other present persons retain \textbf{ēo}. The selected dental combination is \textbf{-tt}; \textbf{bīet} is a documented valid third-singular alternative. [3]
% mdparagraph:P032:end

\end{minipage}\par\addvspace{4pt}

\par\begin{minipage}{\linewidth}\raggedright
% mdparagraph:P033:begin
\textbf{B2.}
% mdparagraph:P033:end

% item: NEW-U04-B02
% key: KEY-NEW-U04-B02
% mdtable:T10:begin
\par\begingroup
\fontsize{10.5}{13.2}\selectfont
\setlength{\LTleft}{0pt}\setlength{\LTright}{\fill}
\noindent
\begin{tabular}{@{}>{\raggedright\arraybackslash}p{21.11mm}>{\raggedright\arraybackslash}p{21.11mm}>{\raggedright\arraybackslash}p{21.11mm}>{\raggedright\arraybackslash}p{22.8mm}>{\raggedright\arraybackslash}p{26.6mm}@{}}
\toprule
\textbf{Present 1sg} & \textbf{Present 2sg} & \textbf{Present 3sg} & \textbf{Preterite 2sg} & \textbf{Preterite 1/\allowbreak{}2/\allowbreak{}3pl} \\
\midrule
būge & bȳgst & bȳgð & buge & bugon \\
\bottomrule
\end{tabular}
\par\endgroup

% mdtable:T10:end
% mdparagraph:P034:begin
Selected \textbf{būgan} has \textbf{ū → ȳ} in the present second and third singular, whereas \textbf{bēodan} has \textbf{ēo → īe}. The preterite plural-based \textbf{u} is short. \textbf{Bȳhst/\allowbreak{}bȳhð} are documented valid spellings beside selected \textbf{bȳgst/\allowbreak{}bȳgð}; the stated environment permits this \textbf{g/\allowbreak{}h} choice. [3]
% mdparagraph:P034:end

\end{minipage}\par\addvspace{4pt}

\par\begin{minipage}{\linewidth}\raggedright
% mdparagraph:P035:begin
\textbf{B3.}
% mdparagraph:P035:end

% item: NEW-U04-B03
% key: KEY-NEW-U04-B03
% mdtable:T11:begin
\par\begingroup
\fontsize{10.5}{13.2}\selectfont
\setlength{\LTleft}{0pt}\setlength{\LTright}{\fill}
\noindent
\begin{tabular}{@{}>{\raggedright\arraybackslash}p{20.8mm}>{\raggedright\arraybackslash}p{20.8mm}>{\raggedright\arraybackslash}p{22.47mm}>{\raggedright\arraybackslash}p{22.47mm}>{\raggedright\arraybackslash}p{26.21mm}@{}}
\toprule
\textbf{Present 2sg} & \textbf{Present 3sg} & \textbf{Preterite 1sg} & \textbf{Preterite 2sg} & \textbf{Preterite 1/\allowbreak{}2/\allowbreak{}3pl} \\
\midrule
cīest & cīest & cēas & cure & curon \\
\bottomrule
\end{tabular}
\par\endgroup

% mdtable:T11:end
% mdparagraph:P036:begin
The present and the first/\allowbreak{}third singular preterite retain \textbf{s}; the plural-preterite stem has \textbf{r}, including second-singular \textbf{cure}. Thus the second and third singular present coincide, but the singular and plural past stems differ. [3]
% mdparagraph:P036:end

\end{minipage}\par\addvspace{4pt}

\par\begin{minipage}{\linewidth}\raggedright
% mdparagraph:P037:begin
\textbf{B4.}
% mdparagraph:P037:end

% item: NEW-U04-B04
% key: KEY-NEW-U04-B04
% mdtable:T12:begin
\par\begingroup
\fontsize{10.5}{13.2}\selectfont
\setlength{\LTleft}{0pt}\setlength{\LTright}{\fill}
\noindent
\begin{tabular}{@{}>{\raggedright\arraybackslash}p{17.12mm}>{\raggedright\arraybackslash}p{17.12mm}>{\raggedright\arraybackslash}p{17.12mm}>{\raggedright\arraybackslash}p{18.49mm}>{\raggedright\arraybackslash}p{18.49mm}>{\raggedright\arraybackslash}p{21.58mm}@{}}
\toprule
\textbf{Verb} & \textbf{Present 2sg} & \textbf{Present 3sg} & \textbf{Preterite 1sg} & \textbf{Preterite 2sg} & \textbf{Preterite 1/\allowbreak{}2/\allowbreak{}3pl} \\
\midrule
sēoðan & sīest & sīeð & sēað & sude & sudon \\
drēogan & drīegst & drīegð & drēag & druge & drugon \\
\bottomrule
\end{tabular}
\par\endgroup

% mdtable:T12:end
% mdparagraph:P038:begin
\textbf{Sēoðan} changes \textbf{ð} to \textbf{d} in its plural-based past forms. Its contracted present endings are separate from that alternation. \textbf{Drēogan} retains \textbf{g} throughout the selected row; \textbf{drīehst, drīehð} and singular preterite \textbf{drēah} are documented alternatives in their stated positions. Those alternatives do not replace every \textbf{g} in the paradigm. Both verbs select long \textbf{īe} in the present second/\allowbreak{}third singular and short \textbf{u} in the plural-based preterite. [3]
% mdparagraph:P038:end

\end{minipage}\par\addvspace{4pt}

\par\begin{minipage}{\linewidth}\raggedright
% mdparagraph:P039:begin
\textbf{B5.}
% mdparagraph:P039:end

% item: NEW-U04-B05
% key: KEY-NEW-U04-B05
% mdtable:T13:begin
% course-table-style: 2026-10
% course-panel:T13-P01:begin
\par\begingroup
\fontsize{10.5}{13.2}\selectfont
\setlength{\LTleft}{0pt}\setlength{\LTright}{\fill}
\noindent
\begin{tabular}{@{}>{\raggedright\arraybackslash}p{39.459mm}>{\raggedright\arraybackslash}p{39.459mm}>{\raggedright\arraybackslash}p{39.459mm}@{}}
\toprule
\textbf{Person and number} & \textbf{Present} & \textbf{Preterite} \\
\midrule
1 sg. & tēo & tēah \\
2 sg. & tīehst & tuge \\
3 sg. & tīehð & tēah \\
1–3 pl. & tēoð & tugon \\
\bottomrule
\end{tabular}
\par\endgroup
% course-panel:T13-P01:end
% mdtable:T13:end
% mdparagraph:P040:begin
These are the Class II forms of \textbf{tēon}, “draw.” Present \textbf{tēo/\allowbreak{}tēoð} are contracted. The preterite contrasts singular \textbf{tēah} with plural-based \textbf{tuge/\allowbreak{}tugon}; it does not retain the present vowel throughout. [3]
% mdparagraph:P040:end

\end{minipage}\par\addvspace{4pt}

\par\begin{minipage}{\linewidth}\raggedright
% mdparagraph:P041:begin
\textbf{B6.}
% mdparagraph:P041:end

% item: NEW-U04-B06
% key: KEY-NEW-U04-B06
% mdtable:T14:begin
\par\begingroup
\fontsize{10.5}{13.2}\selectfont
\setlength{\LTleft}{0pt}\setlength{\LTright}{\fill}
\noindent
\begin{tabular}{@{}>{\raggedright\arraybackslash}p{20.8mm}>{\raggedright\arraybackslash}p{20.8mm}>{\raggedright\arraybackslash}p{22.47mm}>{\raggedright\arraybackslash}p{22.47mm}>{\raggedright\arraybackslash}p{26.21mm}@{}}
\toprule
\textbf{Verb} & \textbf{Present 1sg} & \textbf{Preterite 1sg} & \textbf{Preterite 2sg} & \textbf{Preterite 1/\allowbreak{}2/\allowbreak{}3pl} \\
\midrule
wrēon & wrēo & wrāh & wrige & wrigon \\
tēon, draw & tēo & tēah & tuge & tugon \\
\bottomrule
\end{tabular}
\par\endgroup

% mdtable:T14:end
% mdparagraph:P042:begin
The shared contracted \textbf{-ēon} ending does not identify the strong class. Selected \textbf{wrēon} follows Class I, with singular past \textbf{ā} and plural-based \textbf{i}; \textbf{tēon} follows Class II, with \textbf{ēa} and \textbf{u}. Preserve the complete chosen stem series instead of combining one verb's present with the other's past pattern. [2, 3]
% mdparagraph:P042:end

\end{minipage}\par\addvspace{4pt}

\reserve{7\baselineskip}
% mdparagraph:P043:begin
\section*{Class III}
% mdparagraph:P043:end

\reserve{11\baselineskip}
\par\begin{minipage}{\linewidth}\raggedright
% mdparagraph:P044:begin
\textbf{C1.}
% mdparagraph:P044:end

% item: NEW-U04-C01
% key: KEY-NEW-U04-C01
% mdtable:T15:begin
% course-table-style: 2026-10
% course-panel:T15-P01:begin
\par\begingroup
\fontsize{10.5}{13.2}\selectfont
\setlength{\LTleft}{0pt}\setlength{\LTright}{\fill}
\noindent
\begin{tabular}{@{}>{\raggedright\arraybackslash}p{39.459mm}>{\raggedright\arraybackslash}p{39.459mm}>{\raggedright\arraybackslash}p{39.459mm}@{}}
\toprule
\textbf{Person and number} & \textbf{Present} & \textbf{Preterite} \\
\midrule
1 sg. & binde & band \\
2 sg. & bintst & bunde \\
3 sg. & bint & band \\
1–3 pl. & bindað & bundon \\
\bottomrule
\end{tabular}
\par\endgroup
% course-panel:T15-P01:end
% mdtable:T15:end
% mdparagraph:P045:begin
In selected \textbf{bintst}, stem-final \textbf{d} combines with \textbf{-st} as \textbf{-tst}. In \textbf{bint}, \textbf{d + ð} gives the dental outcome written here with one \textbf{t}. Cook also records \textbf{bindest/\allowbreak{}bindst} and \textbf{bindeð/\allowbreak{}bindð}; these are valid alternatives, though the task requests the selected row. The past second singular \textbf{bunde} uses the \textbf{u} stem, not singular \textbf{band}. [4]
% mdparagraph:P045:end

\end{minipage}\par\addvspace{4pt}

\par\begin{minipage}{\linewidth}\raggedright
% mdparagraph:P046:begin
\textbf{C2.}
% mdparagraph:P046:end

% item: NEW-U04-C02
% key: KEY-NEW-U04-C02
% mdtable:T16:begin
\par\begingroup
\fontsize{10.5}{13.2}\selectfont
\setlength{\LTleft}{0pt}\setlength{\LTright}{\fill}
\noindent
\begin{tabular}{@{}>{\raggedright\arraybackslash}p{17.12mm}>{\raggedright\arraybackslash}p{17.12mm}>{\raggedright\arraybackslash}p{17.12mm}>{\raggedright\arraybackslash}p{18.49mm}>{\raggedright\arraybackslash}p{18.49mm}>{\raggedright\arraybackslash}p{21.58mm}@{}}
\toprule
\textbf{Verb} & \textbf{Present 2sg} & \textbf{Present 3sg} & \textbf{Preterite 1sg} & \textbf{Preterite 2sg} & \textbf{Preterite 1/\allowbreak{}2/\allowbreak{}3pl} \\
\midrule
helpan & hilpst & hilpð & healp & hulpe & hulpon \\
gieldan & gieltst & gielt & geald & gulde & guldon \\
\bottomrule
\end{tabular}
\par\endgroup

% mdtable:T16:end
% mdparagraph:P047:begin
\textbf{Helpan} has present \textbf{e/\allowbreak{}i}, while \textbf{gieldan} has selected \textbf{ie} in these present forms. Both have short \textbf{ea} in the first/\allowbreak{}third singular preterite and short \textbf{u} in the plural-based forms. The dental contractions in \textbf{gieltst/\allowbreak{}gielt} belong to the present; keep \textbf{gulde/\allowbreak{}guldon} for the plural-based past. [4]
% mdparagraph:P047:end

\end{minipage}\par\addvspace{4pt}

\par\begin{minipage}{\linewidth}\raggedright
% mdparagraph:P048:begin
\textbf{C3.}
% mdparagraph:P048:end

% item: NEW-U04-C03
% key: KEY-NEW-U04-C03
% mdtable:T17:begin
\par\begingroup
\fontsize{10.5}{13.2}\selectfont
\setlength{\LTleft}{0pt}\setlength{\LTright}{\fill}
\noindent
\begin{tabular}{@{}>{\raggedright\arraybackslash}p{17.12mm}>{\raggedright\arraybackslash}p{17.12mm}>{\raggedright\arraybackslash}p{17.12mm}>{\raggedright\arraybackslash}p{18.49mm}>{\raggedright\arraybackslash}p{18.49mm}>{\raggedright\arraybackslash}p{21.58mm}@{}}
\toprule
\textbf{Verb} & \textbf{Present 2sg} & \textbf{Present 3sg} & \textbf{Preterite 1sg} & \textbf{Preterite 2sg} & \textbf{Preterite 1/\allowbreak{}2/\allowbreak{}3pl} \\
\midrule
weorpan & wierpst & wierpð & wearp & wurpe & wurpon \\
berstan & bierst & bierst & bærst & burste & burston \\
\bottomrule
\end{tabular}
\par\endgroup

% mdtable:T17:end
% mdparagraph:P049:begin
The first/\allowbreak{}third preterite of \textbf{weorpan} has short \textbf{ea}. The metathesized \textbf{berst-} pattern has short \textbf{æ}, giving \textbf{bærst}. Both use short \textbf{u} for the plural-based preterite. The \textbf{ie} in \textbf{wierpst/\allowbreak{}wierpð/\allowbreak{}bierst} is short, unlike Class II \textbf{īe}. [4]
% mdparagraph:P049:end

\end{minipage}\par\addvspace{4pt}

\par\begin{minipage}{\linewidth}\raggedright
% mdparagraph:P050:begin
\textbf{C4.}
% mdparagraph:P050:end

% item: NEW-U04-C04
% key: KEY-NEW-U04-C04
% mdtable:T18:begin
% course-table-style: 2026-10
% course-panel:T18-P01:begin
\par\begingroup
\fontsize{10.5}{13.2}\selectfont
\setlength{\LTleft}{0pt}\setlength{\LTright}{\fill}
\noindent
\begin{tabular}{@{}>{\raggedright\arraybackslash}p{39.459mm}>{\raggedright\arraybackslash}p{39.459mm}>{\raggedright\arraybackslash}p{39.459mm}@{}}
\toprule
\textbf{Person and number} & \textbf{Present} & \textbf{Preterite} \\
\midrule
1 sg. & weorðe & wearð \\
2 sg. & wierst & wurde \\
3 sg. & wierð & wearð \\
1–3 pl. & weorðað & wurdon \\
\bottomrule
\end{tabular}
\par\endgroup
% course-panel:T18-P01:end
% mdtable:T18:end
% mdparagraph:P051:begin
\textbf{Ð} appears in present \textbf{weorðe, wierð, weorðað} and first/\allowbreak{}third singular preterite \textbf{wearð}. The selected second-singular present \textbf{wierst} has no \textbf{ð}; retained \textbf{wierðst} is also valid. \textbf{D} occurs in preterite \textbf{wurde} and all three plural persons \textbf{wurdon}. The productive sense here is ordinary “become”; these forms do not by themselves teach a passive construction. [4]
% mdparagraph:P051:end

\end{minipage}\par\addvspace{4pt}

\par\begin{minipage}{\linewidth}\raggedright
% mdparagraph:P052:begin
\textbf{C5.}
% mdparagraph:P052:end

% item: NEW-U04-C05
% key: KEY-NEW-U04-C05
% mdtable:T19:begin
% course-table-style: 2026-10
% course-panel:T19-P01:begin
\par\begingroup
\fontsize{10.5}{13.2}\selectfont
\setlength{\LTleft}{0pt}\setlength{\LTright}{\fill}
\noindent
\begin{tabular}{@{}>{\raggedright\arraybackslash}p{39.459mm}>{\raggedright\arraybackslash}p{39.459mm}>{\raggedright\arraybackslash}p{39.459mm}@{}}
\toprule
\textbf{Person and number} & \textbf{Present} & \textbf{Preterite} \\
\midrule
1 sg. & finde & fand \\
2 sg. & fintst & funde \\
3 sg. & fint & funde \\
1–3 pl. & findað & fundon \\
\bottomrule
\end{tabular}
\par\endgroup
% course-panel:T19-P01:end
% mdtable:T19:end
% mdparagraph:P053:begin
The selected usual third-singular preterite is \textbf{funde}; \textbf{fand} is also valid. Selected first singular is \textbf{fand}, but Wright also attests first-singular \textbf{funde}. A response using either documented alternative can be valid Old English while differing from the requested selected table. The second singular is \textbf{funde}, and all plural persons have \textbf{fundon}. [4]
% mdparagraph:P053:end

% mdparagraph:P054:begin
For \textbf{winnan}:
% mdparagraph:P054:end

% mdtable:T20:begin
\par\begingroup
\fontsize{10.5}{13.2}\selectfont
\setlength{\LTleft}{0pt}\setlength{\LTright}{\fill}
\noindent
\begin{tabular}{@{}>{\raggedright\arraybackslash}p{36.31mm}>{\raggedright\arraybackslash}p{36.31mm}>{\raggedright\arraybackslash}p{45.75mm}@{}}
\toprule
\textbf{Present 2sg} & \textbf{Present 3sg} & \textbf{Preterite 1/\allowbreak{}2/\allowbreak{}3pl} \\
\midrule
winst & winð & wunnon \\
\bottomrule
\end{tabular}
\par\endgroup

% mdtable:T20:end
% mdparagraph:P055:begin
The doubled \textbf{nn} of the present stem is reduced before the consonantal singular endings. It remains doubled in \textbf{wunnon}.
% mdparagraph:P055:end

\end{minipage}\par\addvspace{4pt}

\par\begin{minipage}{\linewidth}\raggedright
% mdparagraph:P056:begin
\textbf{C6.}
% mdparagraph:P056:end

% item: NEW-U04-C06
% key: KEY-NEW-U04-C06
% mdtable:T21:begin
\par\begingroup
\fontsize{10.5}{13.2}\selectfont
\setlength{\LTleft}{0pt}\setlength{\LTright}{\fill}
\noindent
\begin{tabular}{@{}>{\raggedright\arraybackslash}p{17.12mm}>{\raggedright\arraybackslash}p{17.12mm}>{\raggedright\arraybackslash}p{17.12mm}>{\raggedright\arraybackslash}p{18.49mm}>{\raggedright\arraybackslash}p{18.49mm}>{\raggedright\arraybackslash}p{21.58mm}@{}}
\toprule
\textbf{Verb} & \textbf{Present 2sg} & \textbf{Present 3sg} & \textbf{Preterite 1sg} & \textbf{Preterite 2sg} & \textbf{Preterite 1/\allowbreak{}2/\allowbreak{}3pl} \\
\midrule
drincan & drincst & drincð & dranc & drunce & druncon \\
singan & singst & singð & sang & sunge & sungon \\
\bottomrule
\end{tabular}
\par\endgroup

% mdtable:T21:end
% mdparagraph:P057:begin
Both belong to the nasal Class III family represented by \textbf{bindan}: present \textbf{i}, first/\allowbreak{}third singular preterite \textbf{a}, plural-based preterite \textbf{u}. The supplied participial stem likewise has \textbf{u}, but no participle recall or syntax is required here. [4]
% mdparagraph:P057:end

\end{minipage}\par\addvspace{4pt}

\par\begin{minipage}{\linewidth}\raggedright
% mdparagraph:P058:begin
\textbf{C7.}
% mdparagraph:P058:end

% item: NEW-U04-C07
% key: KEY-NEW-U04-C07
% mdparagraph:P059:begin
\textbf{Tōberstan}, “break up; go to pieces”:
% mdparagraph:P059:end

% mdtable:T22:begin
% course-table-style: 2026-10
% course-panel:T22-P01:begin
\par\begingroup
\fontsize{10.5}{13.2}\selectfont
\setlength{\LTleft}{0pt}\setlength{\LTright}{\fill}
\noindent
\begin{tabular}{@{}>{\raggedright\arraybackslash}p{39.459mm}>{\raggedright\arraybackslash}p{39.459mm}>{\raggedright\arraybackslash}p{39.459mm}@{}}
\toprule
\textbf{Person and number} & \textbf{Present} & \textbf{Preterite} \\
\midrule
1 sg. & tōberste & tōbærst \\
2 sg. & tōbierst & tōburste \\
3 sg. & tōbierst & tōbærst \\
1–3 pl. & tōberstað & tōburston \\
\bottomrule
\end{tabular}
\par\endgroup
% course-panel:T22-P01:end
% mdtable:T22:end
% mdparagraph:P060:begin
It uses the \textbf{berstan} pattern: present \textbf{berst-/\allowbreak{}bierst-}, past singular \textbf{bærst-}, and plural-based \textbf{burst-}, with its own \textbf{tō-} retained.
% mdparagraph:P060:end

% mdparagraph:P061:begin
\textbf{Forweorðan}, “perish”:
% mdparagraph:P061:end

% mdtable:T23:begin
% course-table-style: 2026-10
% course-panel:T23-P01:begin
\par\begingroup
\fontsize{10.5}{13.2}\selectfont
\setlength{\LTleft}{0pt}\setlength{\LTright}{\fill}
\noindent
\begin{tabular}{@{}>{\raggedright\arraybackslash}p{39.459mm}>{\raggedright\arraybackslash}p{39.459mm}>{\raggedright\arraybackslash}p{39.459mm}@{}}
\toprule
\textbf{Person and number} & \textbf{Present} & \textbf{Preterite} \\
\midrule
1 sg. & forweorðe & forwearð \\
2 sg. & forwierst & forwurde \\
3 sg. & forwierð & forwearð \\
1–3 pl. & forweorðað & forwurdon \\
\bottomrule
\end{tabular}
\par\endgroup
% course-panel:T23-P01:end
% mdtable:T23:end
% mdparagraph:P062:begin
It uses the \textbf{weorðan} pattern, including the \textbf{ð/\allowbreak{}d} alternation. \textbf{Forwierðst} is a valid alternative to selected \textbf{forwierst}. Both whole verbs have separately learned intransitive meanings. Recognizing their base patterns does not authorize attaching prefixes freely or guessing a whole verb's meaning from its parts. [4, 7]
% mdparagraph:P062:end

\end{minipage}\par\addvspace{4pt}

\par\begin{minipage}{\linewidth}\raggedright
% mdparagraph:P063:begin
\textbf{C8.}
% mdparagraph:P063:end

% item: NEW-U04-C08
% key: KEY-NEW-U04-C08
% mdtable:T24:begin
\par\begingroup
\fontsize{10.5}{13.2}\selectfont
\setlength{\LTleft}{0pt}\setlength{\LTright}{\fill}
\noindent
\begin{tabular}{@{}>{\raggedright\arraybackslash}p{32.25mm}>{\raggedright\arraybackslash}p{38.41mm}>{\raggedright\arraybackslash}p{47.72mm}@{}}
\toprule
\textbf{Reference verb} & \textbf{Selected present 3sg} & \textbf{Selected preterite 1/\allowbreak{}2/\allowbreak{}3pl} \\
\midrule
bregdan & brītt & brugdon \\
fēolan & filhð & fulgon \\
\bottomrule
\end{tabular}
\par\endgroup

% mdtable:T24:end
% mdparagraph:P064:begin
\textbf{Brītt} shows loss of \textbf{g}, lengthening in the exceptional present stem, and the dental combination. The selected \textbf{brugdon} retains \textbf{g}. The alternatives \textbf{brīt} and preterite plural \textbf{brūdon} are documented; the latter belongs to the separately noted family with loss of \textbf{g} and lengthening.
% mdparagraph:P064:end

% mdparagraph:P065:begin
\textbf{Fēolan} has a contracted present citation, short \textbf{i} with retained \textbf{h} in \textbf{filhð}, and \textbf{g} in plural preterite \textbf{fulgon}. Cook also gives plural \textbf{fǣlon}. Its complete reference table combines Cook's parts with the dictionary present entries and ordinary finite endings; it is not a verbatim full paradigm printed by Cook. The dictionary's \textbf{þ} in \textbf{filhþ} is explicitly represented by \textbf{ð} here. Cook's supplied participle \textbf{folen} remains distinct from Wright's \textbf{folgen/\allowbreak{}fōlen} alternatives. These are supported reference observations, not an unaided production demand. [4, 8]
% mdparagraph:P065:end

\end{minipage}\par\addvspace{4pt}

\reserve{7\baselineskip}
% mdparagraph:P066:begin
\section*{Nouns}
% mdparagraph:P066:end

\reserve{9\baselineskip}
\par\begin{minipage}{\linewidth}\raggedright
% mdparagraph:P067:begin
\textbf{N1.}
% mdparagraph:P067:end

% item: NEW-U04-N01
% key: KEY-NEW-U04-N01
% mdtable:T25:begin
\par\begingroup
\fontsize{10.5}{13.2}\selectfont
\setlength{\LTleft}{0pt}\setlength{\LTright}{\fill}
\noindent
\begin{tabular}{@{}>{\raggedright\arraybackslash}p{22.33mm}>{\raggedright\arraybackslash}p{22.89mm}>{\raggedright\arraybackslash}p{22.33mm}>{\raggedright\arraybackslash}p{22.89mm}>{\raggedright\arraybackslash}p{22.33mm}@{}}
\toprule
\textbf{Case} & \textbf{scip singular} & \textbf{scip plural} & \textbf{wind singular} & \textbf{wind plural} \\
\midrule
Nominative & scip & scipu & wind & windas \\
Genitive & scipes & scipa & windes & winda \\
Dative & scipe & scipum & winde & windum \\
Accusative & scip & scipu & wind & windas \\
\bottomrule
\end{tabular}
\par\endgroup

% mdtable:T25:end
% mdparagraph:P068:begin
\textbf{Scip} is a short-stem neuter, with nominative/\allowbreak{}accusative plural \textbf{-u}. Its initial \textbf{sc} is an onset; it does not make the syllable heavy. Strong masculine \textbf{wind} has nominative/\allowbreak{}accusative plural \textbf{-as}. Neither noun should take the other's plural ending. [5]
% mdparagraph:P068:end

\end{minipage}\par\addvspace{4pt}

\par\begin{minipage}{\linewidth}\raggedright
% mdparagraph:P069:begin
\textbf{N2.}
% mdparagraph:P069:end

% item: NEW-U04-N02
% key: KEY-NEW-U04-N02
% mdtable:T26:begin
\par\begingroup
\fontsize{10.5}{13.2}\selectfont
\setlength{\LTleft}{0pt}\setlength{\LTright}{\fill}
\noindent
\begin{tabular}{@{}>{\raggedright\arraybackslash}p{21.35mm}>{\raggedright\arraybackslash}p{23.81mm}>{\raggedright\arraybackslash}p{21.89mm}>{\raggedright\arraybackslash}p{23.81mm}>{\raggedright\arraybackslash}p{21.89mm}@{}}
\toprule
\textbf{Case} & \textbf{strand singular} & \textbf{strand plural} & \textbf{gefēra singular} & \textbf{gefēra plural} \\
\midrule
Nominative & strand & strandas & gefēra & gefēran \\
Genitive & strandes & stranda & gefēran & gefērena \\
Dative & strande & strandum & gefēran & gefērum \\
Accusative & strand & strandas & gefēran & gefēran \\
\bottomrule
\end{tabular}
\par\endgroup

% mdtable:T26:end
% mdparagraph:P070:begin
\textbf{Gefēran} can be genitive, dative or accusative singular, or nominative or accusative plural. Its spelling alone therefore fixes neither case nor number. Agreement and construction distinguish the possibilities. Nominative singular \textbf{gefēra}, genitive plural \textbf{gefērena} and dative plural \textbf{gefērum} remain distinct. [5]
% mdparagraph:P070:end

\end{minipage}\par\addvspace{4pt}

\par\begin{minipage}{\linewidth}\raggedright
% mdparagraph:P071:begin
\textbf{N3.}
% mdparagraph:P071:end

% item: NEW-U04-N03
% key: KEY-NEW-U04-N03
% mdparagraph:P072:begin
\textbf{Ðā gefēran fundon twā scipu on ðǣm strande. Ðæs gōdan mannes scip wæs lȳtel.}
% mdparagraph:P072:end

% mdparagraph:P073:begin
The companions are nominative plural; \textbf{twā scipu} is the accusative plural object. \textbf{On ðǣm strande} locates the ships/\allowbreak{}finding on the shore and has dative singular, not a newly introduced destination use. In the second sentence, \textbf{ðæs gōdan mannes} is the complete limiting genitive before \textbf{scip}. \textbf{Lȳtel} is the strong neuter nominative singular predicate agreeing with \textbf{scip}. A postposed genitive could be grammatical but would miss the requested position.
% mdparagraph:P073:end

\end{minipage}\par\addvspace{4pt}

\reserve{7\baselineskip}
% mdparagraph:P074:begin
\section*{Uses and connected sentences}
% mdparagraph:P074:end

\reserve{5\baselineskip}
\par\begin{minipage}{\linewidth}\raggedright
% mdparagraph:P075:begin
\textbf{D1.}
% mdparagraph:P075:end

% item: NEW-U04-D01
% key: KEY-NEW-U04-D01
% mdparagraph:P076:begin
\textbf{Ðā gefēran cēosað ðæt lȳtle hof. Wē cēosað ðæt micle scip.}
% mdparagraph:P076:end

% mdparagraph:P077:begin
Both subjects are plural and require \textbf{cēosað}. The things chosen are accusative neuter singular; the adjectives are weak after \textbf{ðæt}. Preserve \textbf{micle}, with its shortened stem before the ending.
% mdparagraph:P077:end

\end{minipage}\par\addvspace{4pt}

\par\begin{minipage}{\linewidth}\raggedright
% mdparagraph:P078:begin
\textbf{D2.}
% mdparagraph:P078:end

% item: NEW-U04-D02
% key: KEY-NEW-U04-D02
% mdparagraph:P079:begin
\textbf{Se mann funde ðæs gōdan mannes scip. Hē band ðæt dēor.}
% mdparagraph:P079:end

% mdparagraph:P080:begin
\textbf{Se mann} is the finder and binder. \textbf{Ðæs gōdan mannes} identifies the ship's owner and precedes \textbf{scip}, the accusative object. \textbf{Hē} continues the first man's reference; \textbf{ðæt dēor} is the singular accusative object of \textbf{band}. Third-singular \textbf{fand} is a valid alternative to selected \textbf{funde}. Postposing the limiting genitive would miss this task's positioning request even if the sentence remained grammatical.
% mdparagraph:P080:end

\end{minipage}\par\addvspace{4pt}

\par\begin{minipage}{\linewidth}\raggedright
% mdparagraph:P081:begin
\textbf{D3.}
% mdparagraph:P081:end

% item: NEW-U04-D03
% key: KEY-NEW-U04-D03
% mdparagraph:P082:begin
\textbf{Wit druncon ðæt wæter, and git sungon wel.}
% mdparagraph:P082:end

% mdparagraph:P083:begin
\textbf{Wit} means “we two,” and \textbf{git} means “you two.” Both take plural verb forms: \textbf{druncon} is first person and \textbf{sungon} second person in these clauses. \textbf{Ðæt wæter} is accusative singular. Replacing either dual pronoun with an ordinary plural loses a requested distinction.
% mdparagraph:P083:end

\end{minipage}\par\addvspace{4pt}

\par\begin{minipage}{\linewidth}\raggedright
% mdparagraph:P084:begin
\textbf{D4.}
% mdparagraph:P084:end

% item: NEW-U04-D04
% key: KEY-NEW-U04-D04
% mdparagraph:P085:begin
\textbf{Ðā męnn bindað ðæt dēor. Ðā fuglas singað wel.}
% mdparagraph:P085:end

% mdparagraph:P086:begin
Only the subjects become plural: \textbf{męnn} and \textbf{fuglas}, with plural \textbf{bindað/\allowbreak{}singað}. \textbf{Ðæt dēor} remains the singular object. Its noun spelling alone could also occur in the plural, but the unchanged article \textbf{ðæt} marks the intended singular phrase.
% mdparagraph:P086:end

\end{minipage}\par\addvspace{4pt}

\par\begin{minipage}{\linewidth}\raggedright
% mdparagraph:P087:begin
\textbf{D5.}
% mdparagraph:P087:end

% item: NEW-U04-D05
% key: KEY-NEW-U04-D05
% mdparagraph:P088:begin
\textbf{Ic rād wīde and fand ān hof.}
% mdparagraph:P088:end

% mdparagraph:P089:begin
The shared subject \textbf{ic} governs both first-singular preterites. \textbf{Ān hof} remains the accusative object of “found”; \textbf{wīde} modifies the riding. \textbf{Funde} is a documented first-singular alternative to selected \textbf{fand}, so it is not inherently ungrammatical.
% mdparagraph:P089:end

\end{minipage}\par\addvspace{4pt}

\par\begin{minipage}{\linewidth}\raggedright
% mdparagraph:P090:begin
\textbf{D6.}
% mdparagraph:P090:end

% item: NEW-U04-D06
% key: KEY-NEW-U04-D06
% mdparagraph:P091:begin
\textbf{Hē singð wel, and se gefēra hīerð hine.}
% mdparagraph:P091:end

% mdparagraph:P092:begin
\textbf{Hē} is the singer, subject of \textbf{singð}. The different person \textbf{se gefēra} is the listener, subject of \textbf{hīerð}. Accusative singular \textbf{hine} refers to the singer: the companion hears him. Both clauses are present, but \textbf{singan} is strong and \textbf{hīeran} weak.
% mdparagraph:P092:end

\end{minipage}\par\addvspace{4pt}

\par\begin{minipage}{\linewidth}\raggedright
% mdparagraph:P093:begin
\textbf{D7.}
% mdparagraph:P093:end

% item: NEW-U04-D07
% key: KEY-NEW-U04-D07
% mdparagraph:P094:begin
\textbf{Se gefēra wearð glæd. Ðā gefēran weorðað glade.}
% mdparagraph:P094:end

% mdparagraph:P095:begin
The predicates agree with their subjects: \textbf{glæd} is strong masculine nominative singular; \textbf{glade} is strong masculine nominative plural. The subject's article does not make the separate predicate adjective weak. \textbf{Wearð/\allowbreak{}weorðað} mean “became/\allowbreak{}become”; there is no passive participle here. [6]
% mdparagraph:P095:end

\end{minipage}\par\addvspace{4pt}

\par\begin{minipage}{\linewidth}\raggedright
% mdparagraph:P096:begin
\textbf{D8.}
% mdparagraph:P096:end

% item: NEW-U04-D08
% key: KEY-NEW-U04-D08
% mdparagraph:P097:begin
\textbf{Ic healp ðǣm ealdan męnn. Se ealda mann healp mē.}
% mdparagraph:P097:end

% mdparagraph:P098:begin
The person helped is dative in the taught construction: \textbf{ðǣm ealdan męnn}, then \textbf{mē}. In the second clause, \textbf{se ealda mann} is nominative, the helper. \textbf{Healp} serves both first and third singular preterite. Genitive objects of \textbf{helpan} are also attested and can be grammatical, but do not meet this task's requested dative construction. [6]
% mdparagraph:P098:end

\end{minipage}\par\addvspace{4pt}

\par\begin{minipage}{\linewidth}\raggedright
% mdparagraph:P099:begin
\textbf{D9.}
% mdparagraph:P099:end

% item: NEW-U04-D09
% key: KEY-NEW-U04-D09
% mdparagraph:P100:begin
\textbf{Ðā gefēran ne fundon ðæt scip. Ic ne cēose ðæt micle hof.}
% mdparagraph:P100:end

% mdparagraph:P101:begin
\textbf{Ne} immediately precedes each finite verb. The first clause is plural preterite, the second first-singular present. Negation does not change the accusative objects or license a weak dental past for the strong verb.
% mdparagraph:P101:end

\end{minipage}\par\addvspace{4pt}

\par\begin{minipage}{\linewidth}\raggedright
% mdparagraph:P102:begin
\textbf{D10.}
% mdparagraph:P102:end

% item: NEW-U04-D10
% key: KEY-NEW-U04-D10
% mdparagraph:P103:begin
\textbf{Ðū ne drunce ðæt wæter. Gē drincað wæter.}
% mdparagraph:P103:end

% mdparagraph:P104:begin
The first subject is one addressee, with second-singular preterite \textbf{drunce}. The second subject is several addressees, with second-plural present \textbf{drincað}. The first object is definite “the water”; the second is simply “water.”
% mdparagraph:P104:end

\end{minipage}\par\addvspace{4pt}

\reserve{7\baselineskip}
% mdparagraph:P105:begin
\section*{Two companions choose a dwelling}
% mdparagraph:P105:end

\reserve{5\baselineskip}
\par\begin{minipage}{\linewidth}\raggedright
% mdparagraph:P106:begin
\textbf{L1.}
% mdparagraph:P106:end

% item: NEW-U04-L01
% key: KEY-NEW-U04-L01
% mdparagraph:P107:begin
Two companions rode far and wide. They found two dwellings. One dwelling was large, and one dwelling was small. They chose the small dwelling. They had food and water. They drank the water and rested in the small dwelling. One man said: “We two have a good dwelling. I sing, and you hear me.”
% mdparagraph:P107:end

% mdparagraph:P108:begin
The two companions remain the subjects of the riding, finding, choosing, having, drinking and resting. \textbf{Geręston hīe} can be rendered literally “rested themselves”; natural English “rested” preserves the meaning. The quoted \textbf{wit} includes the speaker and the other companion. \textbf{Ic singe} and \textbf{ðū hīerst mē} describe present actions: “I am singing, and you hear me” is also possible. Do not turn them into a command, obligation or future promise. This translation is independently prepared from the new Old English account.
% mdparagraph:P108:end

\end{minipage}\par\addvspace{4pt}

\par\begin{minipage}{\linewidth}\raggedright
% mdparagraph:P109:begin
\textbf{L2.}
% mdparagraph:P109:end

% item: NEW-U04-L02
% key: KEY-NEW-U04-L02
% mdtable:T27:begin
\par\begingroup
\fontsize{10.5}{13.2}\selectfont
\setlength{\LTleft}{0pt}\setlength{\LTright}{\fill}
\noindent
\begin{tabular}{@{}>{\raggedright\arraybackslash}p{21.85mm}>{\raggedright\arraybackslash}p{17.62mm}>{\raggedright\arraybackslash}p{17.62mm}>{\raggedright\arraybackslash}p{17.62mm}>{\raggedright\arraybackslash}p{17.62mm}>{\raggedright\arraybackslash}p{17.62mm}@{}}
\toprule
\textbf{Form in the account} & \textbf{Headword} & \textbf{Class} & \textbf{Person} & \textbf{Number} & \textbf{Tense} \\
\midrule
ridon & rīdan & I & Third & Plural & Preterite \\
fundon & findan & III & Third & Plural & Preterite \\
curon & cēosan & II & Third & Plural & Preterite \\
druncon & drincan & III & Third & Plural & Preterite \\
singe & singan & III & First & Singular & Present \\
\bottomrule
\end{tabular}
\par\endgroup

% mdtable:T27:end
% mdparagraph:P110:begin
All five are indicative. \textbf{Hæfdon, geręston, sægde, habbað} and \textbf{hīerst} belong to previously taught weak systems; \textbf{wæs} is the copula. They are not additional strong I–III forms merely because they occur in this account.
% mdparagraph:P110:end

\end{minipage}\par\addvspace{4pt}

\par\begin{minipage}{\linewidth}\raggedright
% mdparagraph:P111:begin
\textbf{L3.}
% mdparagraph:P111:end

% item: NEW-U04-L03
% key: KEY-NEW-U04-L03
% mdtable:T28:begin
\par\begingroup
\fontsize{10.5}{13.2}\selectfont
\setlength{\LTleft}{0pt}\setlength{\LTright}{\fill}
\noindent
\begin{tabular}{@{}>{\raggedright\arraybackslash}p{33.21mm}>{\raggedright\arraybackslash}p{49.31mm}>{\raggedright\arraybackslash}p{35.86mm}@{}}
\toprule
\textbf{Plural form} & \textbf{Selected preterite 3sg} & \textbf{Preterite 2sg} \\
\midrule
ridon & rād & ride \\
fundon & funde & funde \\
curon & cēas & cure \\
druncon & dranc & drunce \\
\bottomrule
\end{tabular}
\par\endgroup

% mdtable:T28:end
% mdparagraph:P112:begin
The second singular consistently uses the plural-preterite stem with \textbf{-e}. \textbf{Findan} has the selected third-singular exception \textbf{funde}, identical here to the second singular; third-singular \textbf{fand} is also valid but differs from the selected target.
% mdparagraph:P112:end

\end{minipage}\par\addvspace{4pt}

\par\begin{minipage}{\linewidth}\raggedright
% mdparagraph:P113:begin
\textbf{L4.}
% mdparagraph:P113:end

% item: NEW-U04-L04
% key: KEY-NEW-U04-L04
% mdparagraph:P114:begin
They choose the small dwelling, \textbf{ðæt lȳtle hof}. The man who speaks sings; the other companion hears that speaker.
% mdparagraph:P114:end

% mdtable:T29:begin
\par\begingroup
\fontsize{10.5}{13.2}\selectfont
\setlength{\LTleft}{0pt}\setlength{\LTright}{\fill}
\noindent
\begin{tabular}{@{}>{\raggedright\arraybackslash}p{33.48mm}>{\raggedright\arraybackslash}p{87.71mm}@{}}
\toprule
\textbf{Pronoun} & \textbf{Reference in the quotation} \\
\midrule
wit & The speaker and the other companion: “we two” \\
ic & The speaker, subject of “sing” \\
ðū & The other companion, subject of “hear” \\
mē & The speaker as the accusative object of “hear” \\
\bottomrule
\end{tabular}
\par\endgroup

% mdtable:T29:end
% mdparagraph:P115:begin
The first and second persons are determined by the quoted conversation, not by the preceding narrator's use of “they.”
% mdparagraph:P115:end

\end{minipage}\par\addvspace{4pt}

\par\begin{minipage}{\linewidth}\raggedright
% mdparagraph:P116:begin
\textbf{L5.}
% mdparagraph:P116:end

% item: NEW-U04-L05
% key: KEY-NEW-U04-L05
% mdparagraph:P117:begin
\textbf{Twēgen gefēran rīdað wīde. Hīe findað twā hofu. Ān hof is micel, and ān hof is lȳtel. Hīe cēosað ðæt lȳtle hof.}
% mdparagraph:P117:end

% mdparagraph:P118:begin
\textbf{Rīdað, findað} and \textbf{cēosað} are third-person plural present; each \textbf{is} is third-person singular present. The case relations and adjectives remain unchanged. In particular, \textbf{twā hofu} is still accusative plural, and \textbf{ðæt lȳtle hof} remains the definite singular object.
% mdparagraph:P118:end

\end{minipage}\par\addvspace{4pt}

\par\begin{minipage}{\linewidth}\raggedright
% mdparagraph:P119:begin
\textbf{L6.}
% mdparagraph:P119:end

% item: NEW-U04-L06
% key: KEY-NEW-U04-L06
% mdparagraph:P120:begin
\textbf{Wē habbað gōd hof. Wē singað, and gē hīerað ūs.}
% mdparagraph:P120:end

% mdparagraph:P121:begin
\textbf{Wē} denotes the group of speakers, \textbf{gē} the other group addressed, and accusative \textbf{ūs} the speakers as the people heard. The three verbs are plural present. The requested groups call for ordinary plural pronouns, not the earlier two-person \textbf{wit}.
% mdparagraph:P121:end

\end{minipage}\par\addvspace{4pt}

\reserve{7\baselineskip}
% mdparagraph:P122:begin
\section*{The shipwreck and the survivor}
% mdparagraph:P122:end

\reserve{5\baselineskip}
\par\begin{minipage}{\linewidth}\raggedright
% mdparagraph:P123:begin
\textbf{R1.}
% mdparagraph:P123:end

% item: NEW-U04-R01
% key: KEY-NEW-U04-R01
% mdparagraph:P124:begin
In addition to that, northeast winds came, and the fierce southwest wind stood against them, and the whole ship broke apart in this terrible storm. Apollonius's companions all perished, and Apollonius alone came by swimming to Pentapolis, the Cyrenian land, and there he came up on the shore.
% mdparagraph:P124:end

% mdparagraph:P125:begin
“Opposed them” is a natural rendering of \textbf{him ongēan stōd}. The local help guides \textbf{him} toward the northeast winds, but the spelling alone does not force that antecedent or exclude a singular referent. A justified contextual alternative should not be rejected merely because \textbf{him} can also be singular.
% mdparagraph:P125:end

% mdparagraph:P126:begin
\textbf{Forwurdon tō dēaðe} states that the companions died: literal “perished to death” can be compressed to “perished.” In contrast, Apollonius reaches land alive and alone. English “but” may make the contrast explicit where the source uses \textbf{and}. “Swam to Pentapolis” and “came ashore there” are defensible idiomatic renderings of the supplied phrases; they do not change the source's noun \textbf{sunde} or dative \textbf{strande}. Preserve the entire outcome, rather than stopping at the ship's breakup.
% mdparagraph:P126:end

% mdparagraph:P127:begin
This translation is independently prepared from the exact Old English passage. The four supplied later-system preterites and local phrase notes may be consulted; their independent recall is not required. [7]
% mdparagraph:P127:end

\end{minipage}\par\addvspace{4pt}

\par\begin{minipage}{\linewidth}\raggedright
% mdparagraph:P128:begin
\textbf{R2.}
% mdparagraph:P128:end

% item: NEW-U04-R02
% key: KEY-NEW-U04-R02
% mdtable:T30:begin
\par\begingroup
\fontsize{10.5}{13.2}\selectfont
\setlength{\LTleft}{0pt}\setlength{\LTright}{\fill}
\noindent
\begin{tabular}{@{}>{\raggedright\arraybackslash}p{18.3mm}>{\raggedright\arraybackslash}p{20mm}>{\raggedright\arraybackslash}p{9.2mm}>{\raggedright\arraybackslash}p{15mm}>{\raggedright\arraybackslash}p{18mm}>{\raggedright\arraybackslash}p{29.44mm}@{}}
\toprule
\textbf{Form} & \textbf{Whole headword} & \textbf{Class} & \textbf{Person and number} & \textbf{Tense and mood} & \textbf{Subject} \\
\midrule
tōbærst & tōberstan & III & Third singular & Preterite indicative & ðæt scip eall, the whole ship \\
forwurdon & forweorðan & III & Third plural & Preterite indicative & Apollonius gefēran ealle, all Apollonius's companions \\
\bottomrule
\end{tabular}
\par\endgroup

% mdtable:T30:end
% mdparagraph:P129:begin
The ship “broke apart”; the companions “perished.” Neither clause has a direct object. Give the whole prefixed headword, preserving its separately learned meaning.
% mdparagraph:P129:end

\end{minipage}\par\addvspace{4pt}

\par\begin{minipage}{\linewidth}\raggedright
% mdparagraph:P130:begin
\textbf{R3.}
% mdparagraph:P130:end

% item: NEW-U04-R03
% key: KEY-NEW-U04-R03
% mdparagraph:P131:begin
The first \textbf{Apollonius}, before \textbf{gefēran}, is genitive: “Apollonius's companions.” Cook's proper-name note explicitly identifies that case. The second \textbf{Apollonius}, before \textbf{āna becōm}, is nominative: Apollonius is the survivor who reached land.
% mdparagraph:P131:end

% mdparagraph:P132:begin
\textbf{Gefēran} is the nominative plural head of the first subject phrase, with plural \textbf{ealle} and \textbf{forwurdon}. The singular possessor does not determine the verb's number. \textbf{Gefēran} alone could represent several other case/\allowbreak{}number combinations; the clause and agreement identify this one. A foreign name's apparent Latin ending is not a reliable substitute for that analysis. [7]
% mdparagraph:P132:end

\end{minipage}\par\addvspace{4pt}

\par\begin{minipage}{\linewidth}\raggedright
% mdparagraph:P133:begin
\textbf{R4.}
% mdparagraph:P133:end

% item: NEW-U04-R04
% key: KEY-NEW-U04-R04
% mdtable:T31:begin
\par\begingroup
\fontsize{10.5}{13.2}\selectfont
\setlength{\LTleft}{0pt}\setlength{\LTright}{\fill}
\noindent
\begin{tabular}{@{}>{\raggedright\arraybackslash}p{21.98mm}>{\raggedright\arraybackslash}p{31.87mm}>{\raggedright\arraybackslash}p{64.52mm}@{}}
\toprule
\textbf{Noun} & \textbf{Case and number here} & \textbf{Clause function} \\
\midrule
windas & Nominative plural & Subject of cōmon \\
wind & Nominative singular & Head of the subject of stōd \\
scip & Nominative singular & Subject of tōbærst \\
gefēran & Nominative plural & Subject of forwurdon, modified by genitive Apollonius \\
strande & Dative singular & In supplied on ðǣm strande \\
\bottomrule
\end{tabular}
\par\endgroup

% mdtable:T31:end
% mdparagraph:P134:begin
\textbf{Windas} follows \textbf{cōmon} but remains its subject. \textbf{Scip} has identical nominative and accusative singular spelling, and \textbf{gefēran} is also syncretic; construction and agreement resolve their roles. \textbf{Strande} remains dative as printed even when English naturally says “came ashore.”
% mdparagraph:P134:end

\end{minipage}\par\addvspace{4pt}

\par\begin{minipage}{\linewidth}\raggedright
% mdparagraph:P135:begin
\textbf{R5.}
% mdparagraph:P135:end

% item: NEW-U04-R05
% key: KEY-NEW-U04-R05
% mdparagraph:P136:begin
He reaches Pentapolis, the Cyrenian land, \textbf{by swimming}, and the passage leaves him \textbf{on the shore}. The supplied \textbf{mid sunde} expresses the means; \textbf{sunde} is a dative noun, not an infinitive.
% mdparagraph:P136:end

% mdparagraph:P137:begin
\textbf{Āna} means “alone”: Apollonius survives without the companions, who have all perished. It is the supplied weak masculine nominative singular form in this particular sense, not simply a count of “one Apollonius.”
% mdparagraph:P137:end

\end{minipage}\par\addvspace{4pt}

\clearpage
% mdparagraph:P138:begin
\section*{Cumulative review}
% mdparagraph:P138:end

\reserve{8\baselineskip}
\par\begin{minipage}{\linewidth}\raggedright
% mdparagraph:P139:begin
\textbf{X1.}
% mdparagraph:P139:end

% item: NEW-U04-X01
% key: KEY-NEW-U04-X01
% mdtable:T32:begin
\par\begingroup
\fontsize{10.5}{13.2}\selectfont
\setlength{\LTleft}{0pt}\setlength{\LTright}{\fill}
\noindent
\begin{tabular}{@{}>{\raggedright\arraybackslash}p{27.9mm}>{\raggedright\arraybackslash}p{93.29mm}@{}}
\toprule
\textbf{Corrected clause} & \textbf{Diagnosis} \\
\midrule
Ðū drife. & Second-singular preterite needs the plural-based stem plus -e. Drāf is first/\allowbreak{}third singular. \\
Hīe curon ðæt scip. & Third-plural preterite needs curon, with u and r, rather than singular cēas. \\
Se mann bint ðæt dēor. & The singular subject needs third-singular present bint, not plural bindað. \\
\bottomrule
\end{tabular}
\par\endgroup

% mdtable:T32:end
% mdparagraph:P140:begin
The two first clauses stay past, and the third stays present. Documented \textbf{bindeð/\allowbreak{}bindð} are valid alternatives to the selected correction \textbf{bint}.
% mdparagraph:P140:end

\end{minipage}\par\addvspace{4pt}

\par\begin{minipage}{\linewidth}\raggedright
% mdparagraph:P141:begin
\textbf{X2.}
% mdparagraph:P141:end

% item: NEW-U04-X02
% key: KEY-NEW-U04-X02
% mdparagraph:P142:begin
One model:
% mdparagraph:P142:end

% mdparagraph:P143:begin
\textbf{Twēgen gefēran ridon. Hīe fundon twā scipu. Hīe curon ðæt lȳtle scip. Hīe druncon wæter.}
% mdparagraph:P143:end

% mdparagraph:P144:begin
Essential meaning: the same two companions rode, found two ships, chose the small ship, and drank water, in that order. All four actions are narrated with plural preterites. \textbf{Twā scipu} is accusative plural; \textbf{ðæt lȳtle scip} is accusative singular with the weak adjective after the article.
% mdparagraph:P144:end

% mdparagraph:P145:begin
This is an open composition. Repeating the subject noun phrase instead of using \textbf{hīe}, or choosing another clear word order within the taught constructions, is acceptable. Require the four connected past actions, two consistent subjects and the stated objects, not one mandatory string. A singular verb, a present verb, or the choice of both ships would change a requested meaning.
% mdparagraph:P145:end

\end{minipage}\par\addvspace{4pt}

\par\begin{minipage}{\linewidth}\raggedright
% mdparagraph:P146:begin
\textbf{X3.}
% mdparagraph:P146:end

% item: NEW-U04-X03
% key: KEY-NEW-U04-X03
% mdparagraph:P147:begin
\textbf{Se ealda gefēra hilpð mīnum fuglum. Mīne fuglas singað wel. Wē fundon his lȳtel scip on ðǣm strande.}
% mdparagraph:P147:end

% mdparagraph:P148:begin
The old companion is the helper, with singular \textbf{hilpð}. \textbf{Mīnum fuglum} is dative plural, the birds helped; \textbf{mīne fuglas} is nominative plural, the singers. Both possessives mean that the birds belong to the speaker, while their endings agree with the birds' cases and number. \textbf{Wē} identifies the plural finders.
% mdparagraph:P148:end

% mdparagraph:P149:begin
\textbf{His} here identifies a male owner, naturally the old companion in this connected context; it does not agree with the neuter noun \textbf{scip}. An independently established different male referent would also be possible. Selected \textbf{lȳtel} is strong accusative neuter singular; weak \textbf{lȳtle} after the possessive is a valid alternative. \textbf{On ðǣm strande} locates the finding/\allowbreak{}ship on the shore.
% mdparagraph:P149:end

% mdparagraph:P150:begin
The model uses the taught dative \textbf{helpan} construction. A documented genitive construction can be valid Old English, though outside this unit's selected dative practice. The bird application is new teaching material, not a quotation of the dictionary's human example. [6]
% mdparagraph:P150:end

\end{minipage}\par\addvspace{4pt}

\par\begin{minipage}{\linewidth}\raggedright
% mdparagraph:P151:begin
\textbf{X4.}
% mdparagraph:P151:end

% item: NEW-U04-X04
% key: KEY-NEW-U04-X04
% mdtable:T33:begin
\par\begingroup
\fontsize{10.5}{13.2}\selectfont
\setlength{\LTleft}{0pt}\setlength{\LTright}{\fill}
\noindent
\begin{tabular}{@{}>{\raggedright\arraybackslash}p{23.35mm}>{\raggedright\arraybackslash}p{61.87mm}>{\raggedright\arraybackslash}p{33.15mm}@{}}
\toprule
\textbf{Clause} & \textbf{Analysis} & \textbf{Meaning} \\
\midrule
Hē funde. & Third-person singular preterite indicative & He found. \\
Ðū funde. & Second-person singular preterite indicative & You, one person, found. \\
Hīe fundon. & Third-person plural preterite indicative & They found. \\
\bottomrule
\end{tabular}
\par\endgroup

% mdtable:T33:end
% mdparagraph:P152:begin
\textbf{Funde} versus \textbf{fundon} distinguishes singular from plural in these clauses. The identical singular finite form does not distinguish the two persons: \textbf{hē} marks third person and \textbf{ðū} second person. The additional documented first-singular use of \textbf{funde} reinforces the point that the form alone does not identify a unique person. \textbf{Fundon} likewise admits all three plural persons, with \textbf{hīe} selecting the third here.
% mdparagraph:P152:end

\end{minipage}\par\addvspace{4pt}

\clearpage
% mdparagraph:P153:begin
\section*{Source notes}
% mdparagraph:P153:end

\begingroup\fontsize{10}{12.8}\selectfont
% mdparagraph:P154:begin
[1] Cook, \textbf{A First Book in Old English}, second edition (1895), §§92–95, printed pp.54–55 /\allowbreak{} PDF70–71: four stems, finite indicative endings, the plural-based second-singular preterite and the special word-order condition for present plural \textbf{-e}. The ordinary subject-before-verb models here retain \textbf{-að} where prescribed.
% mdparagraph:P154:end

% mdparagraph:P155:begin
[2] Cook §§101–102, printed pp.58–60 /\allowbreak{} PDF74–76, with consonant rules §§28–37, pp.22–25 /\allowbreak{} PDF38–41. Wright and Wright, \textbf{Old English Grammar}, second edition (1914), §§491–492, p.255 /\allowbreak{} PDF275, corroborates the Class I stem relationships and contractions.
% mdparagraph:P155:end

% mdparagraph:P156:begin
[3] Cook §103, printed pp.60–61 /\allowbreak{} PDF76–77. Wright §§493–496, pp.256–257 /\allowbreak{} PDF276–277, records the Class II alternatives. Cook also permits \textbf{ī} for \textbf{īe}; the models keep the selected \textbf{īe} spellings.
% mdparagraph:P156:end

% mdparagraph:P157:begin
[4] Cook §104, printed pp.61–63 /\allowbreak{} PDF77–79. Wright §§497–502, pp.257–259 /\allowbreak{} PDF277–279, supplies the cited exceptional variants, including first-singular \textbf{funde}. A rule-derived finite cell is not a claim that every cell is individually printed in a source.
% mdparagraph:P157:end

% mdparagraph:P158:begin
[5] Cook §§43, 47, 53, printed pp.28, 32, 37 /\allowbreak{} PDF44, 48, 53; dictionary entries \textbf{scip}, p.311 /\allowbreak{} PDF327, \textbf{wind}, p.322 /\allowbreak{} PDF338, \textbf{strand}, p.314 /\allowbreak{} PDF330, and \textbf{ge-fēra}, p.291 /\allowbreak{} PDF307. The full noun tables apply these class patterns to the checked lexical entries.
% mdparagraph:P158:end

% mdparagraph:P159:begin
[6] Cook §150, printed p.88 /\allowbreak{} PDF104, and §168, p.96 /\allowbreak{} PDF112, on predicate agreement and accusative objects. Bosworth–Toller, main dictionary (1898), \textbf{helpan}, p.527b, explicitly admits genitive or dative government; its dative human examples support the selected construction, with birds a declared new application. The other lexical frames are individually sourced in the lesson's evidence, not inferred from membership in a strong class.
% mdparagraph:P159:end

% mdparagraph:P160:begin
[7] Cook, Reader XII, \textbf{Apollonius of Tyre}, “The Shipwreck,” printed pp.165–166 /\allowbreak{} PDF181–182, from p.165 editorial body line6 through p.166 line4, stopping at the first \textbf{on ðǣm strande.} The editorial line numbers expand the printed five-line intervals. Relevant source notes are p.165 n.7 and p.166 nn.1–3; the latter identify the genitive proper name, \textbf{forweorðan}, and §79 for \textbf{āna}. The preceding storm supplies the initial connective's context. Dictionary \textbf{for-weorðan}, p.288 /\allowbreak{} PDF304, and \textbf{tō-berstan}, p.316 /\allowbreak{} PDF332, support the whole prefixed senses.
% mdparagraph:P160:end

% mdparagraph:P161:begin
[8] The \textbf{fēolan} reference combines Cook §104 with Bosworth–Toller main dictionary p.275a (\textbf{felgan}, present \textbf{filhst/\allowbreak{}filhþ}) and Toller Supplement (1921) p.211b, which transfers the cited passages to \textbf{fēolan}. \textbf{Þ → ð} is a declared spelling normalization. The complete table is a synthesis using ordinary endings, not an individually attested twelve-cell paradigm. Cook's \textbf{folen} and Wright's \textbf{folgen/\allowbreak{}fōlen} remain separately identified.
% mdparagraph:P161:end

\endgroup
\end{document}
